\documentclass[12pt]{article}
\usepackage[ukrainian]{babel}
\usepackage{fontspec}
% --- НАЛАШТУВАННЯ ШРИФТІВ ---
\setmainfont{Schoolbook-Regular.otf}[
    Path = ../../,
    BoldFont = Schoolbook-Bold.otf,
    ItalicFont = Schoolbook-Italic.otf,
    BoldItalicFont = Schoolbook-BoldItalic.otf
]
\usepackage[italic]{mathastext}
\usepackage[a4paper,margin=1.8cm,bottom=2cm,top=2cm]{geometry}
\usepackage{amsmath,amssymb}
\usepackage{tikz}
\usepackage{pgfplots}
\pgfplotsset{compat=1.16}
\usetikzlibrary{calc,patterns,angles,quotes}
\usepackage{xcolor,array,fancyhdr,enumitem,multicol}
\usepackage{colortbl}
\usepackage{tcolorbox}
\tcbuselibrary{skins}
\usepackage[hidelinks]{hyperref}
\usepackage[firstpage=false, text=@pvtr2525, scale=0.6, color=gray!12]{draftwatermark}
\definecolor{mainGreen}{RGB}{34, 120, 64}
\definecolor{yearOrange}{RGB}{220, 100, 30}
\definecolor{yearcolor}{RGB}{220, 100, 30}
\definecolor{titleBg}{RGB}{225, 240, 220}
\definecolor{instrBg}{RGB}{255, 240, 220}
\definecolor{instrBorder}{RGB}{220, 150, 80}
\pagestyle{fancy}
\fancyhf{}
\renewcommand{\headrulewidth}{0.4pt}
\renewcommand{\footrulewidth}{0.4pt}
\fancyhead[C]{\small\color{gray!80} Збірник НМТ \textendash{} сесії 23.05--5.06.2026}
\fancyhead[R]{\small\color{mainGreen}\textbf{\href{https://t.me/pvtr2525}{@pvtr2525}}}
\fancyfoot[L]{\small\color{gray!80}\href{https://t.me/pvtr2525}{ТГ: @pvtr2525}}
\fancyfoot[C]{\small\color{gray!80}\thepage}
\fancyfoot[R]{\small\color{gray!80}НМТ 2026}
% --- ТАБЛИЦІ ВІДПОВІДЕЙ ---
\newcommand{\answerTable}[5]{%
\par\vspace{0.15cm}
\noindent
\begin{tabular}{|*{5}{>{\centering\arraybackslash}m{3cm}|}}
\hline
\rule[-0.15cm]{0pt}{0.6cm}\textbf{А} & \rule[-0.15cm]{0pt}{0.6cm}\textbf{Б} & \rule[-0.15cm]{0pt}{0.6cm}\textbf{В} & \rule[-0.15cm]{0pt}{0.6cm}\textbf{Г} & \rule[-0.15cm]{0pt}{0.6cm}\textbf{Д} \\
\hline
\rule[-0.2cm]{0pt}{0.75cm}#1 & \rule[-0.2cm]{0pt}{0.75cm}#2 & \rule[-0.2cm]{0pt}{0.75cm}#3 & \rule[-0.2cm]{0pt}{0.75cm}#4 & \rule[-0.2cm]{0pt}{0.75cm}#5 \\
\hline
\end{tabular}
\par\vspace{0.25cm}
}
\newcommand{\answerTableTall}[5]{%
\par\vspace{0.15cm}
\noindent
\begin{tabular}{|*{5}{>{\centering\arraybackslash}m{3cm}|}}
\hline
\rule[-0.2cm]{0pt}{0.7cm}\textbf{А} & \rule[-0.2cm]{0pt}{0.7cm}\textbf{Б} & \rule[-0.2cm]{0pt}{0.7cm}\textbf{В} & \rule[-0.2cm]{0pt}{0.7cm}\textbf{Г} & \rule[-0.2cm]{0pt}{0.7cm}\textbf{Д} \\
\hline
\rule[-0.45cm]{0pt}{1.2cm}#1 & \rule[-0.45cm]{0pt}{1.2cm}#2 & \rule[-0.45cm]{0pt}{1.2cm}#3 & \rule[-0.45cm]{0pt}{1.2cm}#4 & \rule[-0.45cm]{0pt}{1.2cm}#5 \\
\hline
\end{tabular}
\par\vspace{0.25cm}
}
\newcommand{\instructionBox}[1]{%
\par\vspace{0.3cm}
\begin{tcolorbox}[colback=instrBg, colframe=instrBorder, boxrule=0.6pt, arc=2pt,
    left=8pt, right=8pt, top=4pt, bottom=4pt]
\centering\small\bfseries #1
\end{tcolorbox}
\par\vspace{0.15cm}
}
\newcommand{\sectionTitle}[1]{%
\par\vspace{0.4cm}
\begin{tcolorbox}[colback=titleBg, colframe=titleBg, boxrule=0pt, arc=8pt,
    halign=center, fontupper=\Large\bfseries\color{mainGreen}]
#1
\end{tcolorbox}
\par\vspace{0.3cm}
}
\newcommand{\task}[2]{\noindent\textbf{#1.}\ \ #2\par\vspace{0.2cm}}
% --- НМТ-СТИЛЬ ПОЛЕ ВІДПОВІДІ ---
\newcommand{\nmtAnswerBox}{%
\par\vspace{0.2cm}\noindent
Відповідь:\ \
\begin{tabular}{@{}*{4}{|>{\centering\arraybackslash}p{0.8cm}}|@{\hspace{0.1cm}\textbf{,}\hspace{0.1cm}}*{3}{|>{\centering\arraybackslash}p{0.8cm}}|@{}}
\hline
\rule[-0.2cm]{0pt}{0.7cm} & & & & & & \\
\hline
\end{tabular}
\par\vspace{0.3cm}
}
% --- СІТКА ДЛЯ МАТЧИНГУ ---
\newcommand{\matchingGrid}{%
    \begingroup
    \renewcommand{\arraystretch}{1.15}
    \setlength{\tabcolsep}{4pt}
    \footnotesize
    \begin{tabular}{r|c|c|c|c|c|}
         \multicolumn{1}{c}{} & \multicolumn{1}{c}{\textbf{А}} & \multicolumn{1}{c}{\textbf{Б}} & \multicolumn{1}{c}{\textbf{В}} & \multicolumn{1}{c}{\textbf{Г}} & \multicolumn{1}{c}{\textbf{Д}} \\ \cline{2-6}
         \textbf{1} & \phantom{X} & \phantom{X} & \phantom{X} & \phantom{X} & \phantom{X} \\ \cline{2-6}
         \textbf{2} & & & & & \\ \cline{2-6}
         \textbf{3} & & & & & \\ \cline{2-6}
    \end{tabular}
    \endgroup
}
% --- ДОДАТКОВІ МАКРОСИ ЗБІРНИКА ---
\newcommand{\taskBlock}[1]{%
\par\vspace{0.45cm}
\noindent{\color{mainGreen}\rule{\linewidth}{0.8pt}}\par\vspace{0.12cm}
\noindent{\small\bfseries\color{mainGreen}#1}\par\vspace{0.25cm}
}
\newcounter{zad}
\newcommand{\zadtask}[1]{\stepcounter{zad}\noindent\textbf{\thezad.}\ \ #1\par\vspace{0.2cm}}
\newcommand{\zadnum}{\stepcounter{zad}\textbf{\thezad.}\ \ }
\newcounter{ans}
\newcommand{\ansitem}[1]{\stepcounter{ans}\noindent\textbf{\theans.}\ #1\par\vspace{0.07cm}}
\newcommand{\nmtyear}[1]{\hfill{\small\color{yearOrange}(НМТ #1)}}% --- ДОДАТКОВІ МАКРОСИ (розв'язки, зміст) ---
\tcbuselibrary{breakable}
\newcommand{\solution}[1]{%
\par\vspace{0.12cm}
\begin{tcolorbox}[colback=titleBg!45, colframe=mainGreen!55, boxrule=0.4pt, arc=2pt,
    left=6pt, right=6pt, top=3pt, bottom=3pt, breakable]
\footnotesize\textbf{Короткий розв'язок (оригінал).}\ #1
\end{tcolorbox}
\par\vspace{0.12cm}
}
\setcounter{tocdepth}{2}
\newcommand{\chapterTitle}[1]{%
\clearpage\phantomsection\addcontentsline{toc}{section}{#1}%
\sectionTitle{#1}}
\newcommand{\typeTitle}[1]{%
\phantomsection\addcontentsline{toc}{subsection}{\quad #1}%
\par\vspace{0.3cm}
\noindent{\large\bfseries\color{yearOrange}#1}\par\vspace{0.08cm}
\noindent{\color{yearOrange}\rule{\linewidth}{0.4pt}}\par\vspace{0.2cm}}
\newcommand{\ansTheme}[1]{\par\vspace{0.25cm}\noindent{\bfseries\color{mainGreen}#1}\par\vspace{0.12cm}}
\newcommand{\ansType}[1]{\par\vspace{0.1cm}\noindent{\itshape\color{yearOrange}#1}\par\vspace{0.08cm}}
\begin{document}
\thispagestyle{empty}
\vspace*{1.2cm}
\begin{center}
{\Large\bfseries\color{mainGreen} ЗБІРНИК НМТ-2026 з математики}\\[0.4cm]
{\Huge\bfseries\color{mainGreen} РОЗДІЛ 1. ЧИСЛА І ВИРАЗИ}\\[0.6cm]
{\large Оригінали сесій 23.05--5.06.2026 з короткими розв'язками та авторськими аналогами}\\[1cm]
{\large\color{mainGreen}\textbf{\href{https://t.me/pvtr2525}{Телеграм: @pvtr2525}}}
\end{center}
\clearpage
\phantomsection\addcontentsline{toc}{section}{Зміст}
\sectionTitle{ЗМІСТ}
{\hypersetup{linkcolor=black}\renewcommand{\contentsname}{}\vspace{-1cm}\tableofcontents}
\chapterTitle{РОЗДІЛ 1. ЧИСЛА І ВИРАЗИ}
\typeTitle{Завдання з вибором однієї відповіді (А--Д)}
\taskBlock{Спрощення раціонального виразу --- сесія 23.05, завдання №7}
\zadtask{Спростіть вираз $\dfrac{(2x - 3)^2 - 9}{x}$. \nmtyear{2026}}
\answerTable{$4x - 12$}{$4x - 6$}{$4x$}{$2x - 6$}{$2x - 12$}
\solution{За формулою $(a-b)^2$: $(2x-3)^2-9=4x^2-12x+9-9=4x^2-12x=4x(x-3)$. Поділивши на $x$, дістаємо $4x-12$. Відповідь: \textbf{А}.}
\zadtask{Спростіть вираз $\dfrac{(3x - 2)^2 - 4}{x}$.}
\answerTable{$9x - 4$}{$9x - 12$}{$3x - 4$}{$9x$}{$3x - 12$}
\zadtask{Спростіть вираз $\dfrac{9 - (3 - 2x)^2}{x}$.}
\answerTable{$4x - 12$}{$12 - 4x$}{$12 + 4x$}{$-4x$}{$12 - 2x$}
\taskBlock{Обчислення логарифма --- сесія 23.05, завдання №9}
\zadtask{Обчисліть $\log_{0{,}5} 8$. \nmtyear{2026}}
\answerTable{$-2$}{$-4$}{$-3$}{$\dfrac{16}{3}$}{$3$}
\solution{Оскільки $0{,}5=2^{-1}$, а $8=2^3$, то $\log_{0{,}5}8=\dfrac{\log_2 8}{\log_2 0{,}5}=\dfrac{3}{-1}=-3$. Відповідь: \textbf{В}.}
\zadtask{Обчисліть $\log_{0{,}2} 25$.}
\answerTable{$-2$}{$2$}{$-5$}{$5$}{$\dfrac{1}{2}$}
\zadtask{Обчисліть $\log_{3} \dfrac{1}{81}$.}
\answerTable{$-4$}{$4$}{$-3$}{$\dfrac{1}{4}$}{$3$}
\taskBlock{Ділення з остачею (порції) --- сесія 30.05, завдання №5}
\zadtask{Для приготування однієї порції коктейлю потрібно $40$ г морозива. Яку \textit{найбільшу} кількість таких порцій коктейлю можна приготувати з $0{,}5$ кг морозива?}
\answerTable{$10$}{$11$}{$12$}{$13$}{$20$}
\solution{$0{,}5$ кг $=500$ г. Кількість порцій: $500:40=12{,}5$, тобто ціла частина $12$. Відповідь: \textbf{В} ($12$).}
\zadtask{Для приготування однієї порції смузі потрібно $60$ г ягід. Яку \textit{найбільшу} кількість таких порцій смузі можна приготувати з $0{,}8$ кг ягід?}
\answerTable{$12$}{$13$}{$14$}{$15$}{$20$}
\zadtask{Для випікання одного кексу потрібно $35$ г родзинок. Пекар випік \textit{найбільшу} можливу кількість таких кексів, маючи $0{,}6$ кг родзинок. Скільки \textit{грамів} родзинок у нього залишилося?}
\answerTable{$0$}{$5$}{$10$}{$15$}{$25$}
\taskBlock{Спрощення раціонального виразу --- сесія 30.05, завдання №6}
\zadtask{Спростіть вираз $\dfrac{9 - x^2}{x^2 + 6x + 9}$.}
\answerTableTall{$\dfrac{3-x}{x+3}$}{$\dfrac{x-3}{x+3}$}{$\dfrac{3+x}{x-3}$}{$-\dfrac{x+3}{x-3}$}{$\dfrac{1}{x+3}$}
\solution{Чисельник $9-x^2=(3-x)(3+x)$, знаменник $x^2+6x+9=(x+3)^2$. Тоді $\dfrac{(3-x)(3+x)}{(x+3)^2}=\dfrac{3-x}{x+3}$. Відповідь: \textbf{А}.}
\zadtask{Спростіть вираз $\dfrac{16 - x^2}{x^2 + 8x + 16}$.}
\answerTableTall{$\dfrac{x-4}{x+4}$}{$\dfrac{4-x}{x+4}$}{$\dfrac{4+x}{x-4}$}{$-\dfrac{x+4}{x-4}$}{$\dfrac{1}{x+4}$}
\zadtask{Спростіть вираз $\dfrac{x^2 - 10x + 25}{x^2 - 25}$.}
\answerTableTall{$\dfrac{x-5}{x+5}$}{$\dfrac{x+5}{x-5}$}{$\dfrac{5-x}{x+5}$}{$-\dfrac{x+5}{x-5}$}{$\dfrac{1}{x-5}$}
\taskBlock{Спрощення тригонометричного виразу --- сесія 30.05, завдання №13}
\zadtask{Спростіть вираз $\mathrm{tg}^2 x \cdot \cos^2 x - 1$.}
\answerTableTall{$-\sin^2 x$}{$-\cos^2 x$}{$-1$}{$\cos^2 x$}{$\sin^2 x$}
\solution{$\mathrm{tg}^2 x\cdot\cos^2 x=\dfrac{\sin^2 x}{\cos^2 x}\cdot\cos^2 x=\sin^2 x$. Тоді $\sin^2 x-1=-(1-\sin^2 x)=-\cos^2 x$. Відповідь: \textbf{Б}.}
\zadtask{Спростіть вираз $\mathrm{ctg}^2 x \cdot \sin^2 x - 1$.}
\answerTableTall{$-\sin^2 x$}{$-\cos^2 x$}{$-1$}{$\cos^2 x$}{$\sin^2 x$}
\zadtask{Спростіть вираз $(1 - \cos x)(1 + \cos x) - 1$.}
\answerTableTall{$-\sin^2 x$}{$-\cos^2 x$}{$\sin^2 x$}{$\cos^2 x$}{$-1$}
\taskBlock{Спрощення виразу зі степенями --- сесія 1.06, завдання №4}
\zadtask{Спростіть вираз $\dfrac{6x^3 + 3x^3}{3}$.}
\answerTable{$3x^6$}{$6x^6$}{$5x^3$}{$3x^3$}{$7x^3$}
\solution{$\dfrac{6x^3+3x^3}{3}=\dfrac{9x^3}{3}=3x^3$. Відповідь: \textbf{Г}.}
\zadtask{Спростіть вираз $\dfrac{10x^5 + 5x^5}{5}$.}
\answerTable{$3x^{10}$}{$15x^5$}{$3x^5$}{$5x^5$}{$2x^5$}
\zadtask{Спростіть вираз $\dfrac{4x^2 \cdot 6x^3}{8x}$.}
\answerTable{$3x^4$}{$3x^5$}{$2x^4$}{$24x^5$}{$3x^6$}
\taskBlock{Сюжетна задача на спільну роботу --- сесія 1.06, завдання №11}
\zadtask{У цеху на лінії виробництва пляшкової води працюють два автомати для клейового маркування. Перший маркує $14\,000$ пляшок за годину, а другий~--- $16\,000$ пляшок за годину. Обидва автомати одночасно працювали $12$ хвилин. Скільки всього пляшок було промарковано за цей час?}
\answerTable{$2000$}{$3000$}{$5000$}{$6000$}{$7500$}
\solution{За годину разом маркують $14\,000+16\,000=30\,000$ пляшок. $12$ хв $=\frac{12}{60}=0{,}2$ год, тож $30\,000\cdot 0{,}2=6000$. Відповідь: \textbf{Г}.}
\zadtask{На пакувальній лінії працюють два автомати. Перший пакує $12\,000$ коробок за годину, а другий~--- $18\,000$ коробок за годину. Обидва автомати одночасно працювали $18$ хвилин. Скільки всього коробок було запаковано за цей час?}
\answerTable{$5400$}{$7500$}{$9000$}{$10\,000$}{$30\,000$}
\zadtask{У друкарні працюють два верстати. Перший друкує $9000$ аркушів за годину, а другий~--- $15\,000$ аркушів за годину. Обидва верстати одночасно працювали $20$ хвилин. На скільки \textit{більше} аркушів надрукував за цей час другий верстат, ніж перший?}
\answerTable{$1500$}{$2000$}{$3000$}{$4000$}{$6000$}
\taskBlock{Знаки тригонометричних функцій --- сесія 1.06, завдання №13}
\zadtask{Відомо, що $\cos 210^\circ \cdot \sin\alpha > 0$. Якого з наведених значень може набувати кут $\alpha$? \nmtyear{2026}}
\answerTable{$0^\circ$}{$30^\circ$}{$90^\circ$}{$150^\circ$}{$300^\circ$}
\solution{$\cos 210^\circ<0$, тому для $\cos 210^\circ\cdot\sin\alpha>0$ потрібно $\sin\alpha<0$. Серед варіантів лише $\sin 300^\circ<0$. Відповідь: \textbf{Д}.}
\zadtask{Відомо, що $\cos 120^\circ \cdot \sin\alpha > 0$. Якого з наведених значень може набувати кут $\alpha$?}
\answerTable{$0^\circ$}{$60^\circ$}{$90^\circ$}{$120^\circ$}{$210^\circ$}
\zadtask{Відомо, що $\cos 300^\circ \cdot \sin\alpha < 0$. Якого з наведених значень може набувати кут $\alpha$?}
\answerTable{$30^\circ$}{$90^\circ$}{$150^\circ$}{$180^\circ$}{$210^\circ$}
\taskBlock{Розкриття дужок і спрощення виразу --- сесія 2.06, завдання №2}
\zadtask{Спростіть вираз $15 - 4(2x^2 + 3)$.}
\answerTable{$3 - 8x^2$}{$18 - 8x^2$}{$27 - 8x^2$}{$22x^2 + 33$}{$22x^2 + 3$}
\solution{Розкриваємо дужки: $15-4(2x^2+3)=15-8x^2-12=3-8x^2$. Відповідь: \textbf{А}.}
\zadtask{Спростіть вираз $12 - 5(3x^2 + 2)$.}
\answerTable{$22 - 15x^2$}{$2 - 15x^2$}{$10 - 15x^2$}{$2 + 15x^2$}{$14 - 15x^2$}
\zadtask{Спростіть вираз $3(4x^2 + 5) - 20$.}
\answerTable{$12x^2 - 5$}{$12x^2 + 5$}{$12x^2 - 35$}{$5 - 12x^2$}{$12x^2 - 15$}
\taskBlock{Спрощення раціонального виразу --- сесія 2.06, завдання №15}
\zadtask{Спростіть вираз $\dfrac{x^2 - 4xy + 4y^2}{x - 2y} : (2y - x)$.}
\answerTable{$-1$}{$-x^2 + 4xy - y^2$}{$2y - x$}{$1$}{$x^2 - 4xy + y^2$}
\solution{Чисельник $x^2-4xy+4y^2=(x-2y)^2$, тож дріб дорівнює $x-2y$. Ділення на $(2y-x)=-(x-2y)$ дає $\dfrac{x-2y}{-(x-2y)}=-1$. Відповідь: \textbf{А}.}
\zadtask{Спростіть вираз $\dfrac{x^2 + 6xy + 9y^2}{x + 3y} : (x + 3y)$.}
\answerTable{$-1$}{$1$}{$x + 3y$}{$-(x + 3y)$}{$(x + 3y)^2$}
\zadtask{Спростіть вираз $\dfrac{x - 5y}{x^2 - 10xy + 25y^2} \cdot (5y - x)$.}
\answerTable{$1$}{$x - 5y$}{$-1$}{$5y - x$}{$\dfrac{1}{x - 5y}$}
\taskBlock{Пропорційні величини (швидкість відтворення) --- сесія 3.06, завдання №3}
\zadtask{Вадим планує переглянути відеоролик тривалістю $60$ хвилин. Перед переглядом він збільшив швидкість відтворення в $1{,}5$ раза. Скільки хвилин тепер триватиме перегляд?}
\answerTable{$40$}{$90$}{$45$}{$30$}{$35$}
\solution{Збільшення швидкості в $1{,}5$ раза зменшує час у стільки ж разів: $60 : 1{,}5 = 40$ (хв). Відповідь: \textbf{А} ($40$).}
\zadtask{Оксана планує переглянути відеолекцію тривалістю $80$ хвилин. Перед переглядом вона збільшила швидкість відтворення в $1{,}25$ раза. Скільки хвилин тепер триватиме перегляд?}
\answerTable{$64$}{$100$}{$60$}{$55$}{$75$}
\zadtask{Тарас слухає аудіокнигу тривалістю $90$ хвилин. Перед прослуховуванням він \textit{зменшив} швидкість відтворення до $0{,}75$ від звичайної. Скільки хвилин тепер триватиме прослуховування?}
\answerTable{$67{,}5$}{$105$}{$162$}{$120$}{$115$}
\taskBlock{Дії з мішаними дробами --- сесія 3.06, завдання №7}
\zadtask{Обчисліть $5\dfrac{1}{3} : 1\dfrac{3}{5}$.}
\answerTableTall{$\dfrac{3}{10}$}{$\dfrac{128}{15}$}{$\dfrac{10}{3}$}{$\dfrac{5}{3}$}{$\dfrac{15}{128}$}
\solution{$5\dfrac13 : 1\dfrac35 = \dfrac{16}{3} : \dfrac{8}{5} = \dfrac{16}{3}\cdot\dfrac{5}{8} = \dfrac{10}{3}$. Відповідь: \textbf{В} ($\dfrac{10}{3}$).}
\zadtask{Обчисліть $4\dfrac{2}{3} : 1\dfrac{1}{6}$.}
\answerTableTall{$\dfrac{1}{4}$}{$4$}{$\dfrac{49}{9}$}{$\dfrac{9}{49}$}{$\dfrac{7}{2}$}
\zadtask{Обчисліть $3\dfrac{3}{4} \cdot 2\dfrac{2}{5}$.}
\answerTableTall{$\dfrac{75}{8}$}{$9$}{$\dfrac{25}{16}$}{$6$}{$\dfrac{15}{2}$}
\taskBlock{Обчислення тригонометричного виразу --- сесія 3.06, завдання №13}
\zadtask{Обчисліть значення виразу $2\sin\left(-\dfrac{\pi}{2}\right) + \cos 3\pi$.}
\answerTable{$-3$}{$-2$}{$-1$}{$2$}{$3$}
\solution{$\sin\left(-\dfrac{\pi}{2}\right) = -1$, $\cos 3\pi = -1$, тому $2\cdot(-1) + (-1) = -3$. Відповідь: \textbf{А} ($-3$).}
\zadtask{Обчисліть значення виразу $4\sin\dfrac{3\pi}{2} + \cos 2\pi$.}
\answerTable{$-5$}{$-3$}{$-4$}{$3$}{$5$}
\zadtask{Обчисліть значення виразу $\operatorname{tg}\dfrac{3\pi}{4} + 2\sin\dfrac{\pi}{2}$.}
\answerTable{$-3$}{$1$}{$-1$}{$3$}{$0$}
\taskBlock{Стандартний вигляд числа --- сесія 4.06, завдання №1}
\zadtask{В астрономії для визначення відстані використовують позасистемну одиницю довжини~--- астрономічна одиниця (а.\,о.), що дорівнює середній відстані від Землі до Сонця: $1$ а.\,о. $\approx 149{,}6$~млн~\textit{км}. Запишіть це значення (у \textit{км}) у стандартному вигляді.}
\answerTable{$149{,}6 \cdot 10^6$}{$1{,}496 \cdot 10^8$}{$0{,}1496 \cdot 10^9$}{$1496 \cdot 10^5$}{$14{,}96 \cdot 10^7$}
\solution{$149{,}6$~млн $=149\,600\,000=1{,}496\cdot10^8$. У стандартному вигляді множник має задовольняти $1\leqslant a<10$, тож $1{,}496\cdot10^8$. Відповідь: \textbf{Б}.}
\zadtask{В астрономії використовують одиницю довжини парсек (пк), що приблизно дорівнює $30{,}86$~трлн~\textit{км}. Запишіть це значення (у \textit{км}) у стандартному вигляді.}
\answerTable{$30{,}86 \cdot 10^{12}$}{$3{,}086 \cdot 10^{13}$}{$0{,}3086 \cdot 10^{14}$}{$3086 \cdot 10^{10}$}{$3{,}086 \cdot 10^{11}$}
\zadtask{Діаметр еритроцита людини приблизно дорівнює $0{,}0000075$~\textit{м}. Запишіть це значення (у \textit{м}) у стандартному вигляді.}
\answerTable{$7{,}5 \cdot 10^{-6}$}{$7{,}5 \cdot 10^{6}$}{$0{,}75 \cdot 10^{-5}$}{$75 \cdot 10^{-7}$}{$7{,}5 \cdot 10^{-5}$}
\taskBlock{Спрощення раціонального виразу --- сесія 4.06, завдання №6}
\zadtask{$\dfrac{a^2 - 25b^2}{10b + 2a} =$}
\answerTable{$\dfrac{a - 5b}{2}$}{$2a - 10b$}{$\dfrac{a + 5b}{2}$}{$2a + 10b$}{$\dfrac{a - 5b}{4}$}
\solution{Чисельник $a^2-25b^2=(a-5b)(a+5b)$, знаменник $10b+2a=2(a+5b)$. Скорочуємо на $(a+5b)$: $\dfrac{(a-5b)(a+5b)}{2(a+5b)}=\dfrac{a-5b}{2}$. Відповідь: \textbf{А}.}
\zadtask{$\dfrac{a^2 - 9b^2}{6b + 2a} =$}
\answerTable{$2a - 6b$}{$\dfrac{a + 3b}{2}$}{$\dfrac{a - 3b}{2}$}{$2a + 6b$}{$\dfrac{a - 3b}{4}$}
\zadtask{$\dfrac{a^2 - 16b^2}{4b - a} =$}
\answerTable{$a + 4b$}{$\dfrac{a + 4b}{4}$}{$-a - 4b$}{$a - 4b$}{$\dfrac{a - 4b}{4}$}
\taskBlock{Тригонометричний вираз (формула подвійного кута) --- сесія 4.06, завдання №13}
\zadtask{Якому проміжку належить значення виразу $6 \sin 15^\circ \cos 15^\circ$?}
\answerTable{$(0;\, 1)$}{$[1;\, 2)$}{$[2;\, 3)$}{$[3;\, 5)$}{$[5;\, 8)$}
\solution{За формулою $2\sin\alpha\cos\alpha=\sin2\alpha$: $6\sin15^\circ\cos15^\circ=3\sin30^\circ=3\cdot\dfrac{1}{2}=1{,}5\in[1;\,2)$. Відповідь: \textbf{Б}.}
\zadtask{Якому проміжку належить значення виразу $8 \sin 15^\circ \cos 15^\circ$?}
\answerTable{$(0;\, 1)$}{$[1;\, 2)$}{$[2;\, 3)$}{$[3;\, 5)$}{$[5;\, 8)$}
\zadtask{Якому проміжку належить значення виразу $6\left(\cos^2 15^\circ - \sin^2 15^\circ\right)$?}
\answerTable{$(0;\, 1)$}{$[1;\, 2)$}{$[2;\, 3)$}{$[3;\, 5)$}{$[5;\, 8)$}
\taskBlock{Пропорційне обчислення (швидкість роботи) --- сесія 5.06, завдання №1}
\zadtask{Настінний принтер друкує зображення зі сталою швидкістю $3$ м$^2$ за одну годину. Скільки \textit{хвилин} принтер друкував на стіні будинку зображення площею $2{,}4$ м$^2$?}
\answerTable{$36$}{$40$}{$51$}{$48$}{$45$}
\solution{Час друку: $\dfrac{2{,}4}{3}=0{,}8$ год $=0{,}8\cdot 60=48$ хв. Відповідь: \textbf{Г}.}
\zadtask{Насос перекачує воду зі сталою швидкістю $5$ м$^3$ за одну годину. Скільки \textit{хвилин} насос наповнював резервуар об'ємом $3{,}5$ м$^3$?}
\answerTable{$35$}{$40$}{$42$}{$45$}{$48$}
\zadtask{Верстат обробляє поверхню зі сталою швидкістю. За $36$ \textit{хвилин} він обробив панель площею $2{,}4$ м$^2$. Скільки квадратних метрів поверхні верстат обробляє за одну \textit{годину}?}
\answerTable{$3{,}6$}{$4$}{$4{,}8$}{$4{,}4$}{$5$}
\taskBlock{Спрощення виразу (формули скороченого множення) --- сесія 5.06, завдання №6}
\zadtask{Спростіть вираз $(a+6)(a-6) - (a-4)^2$.}
\answerTable{$8a - 52$}{$-8a - 52$}{$8a - 20$}{$-52$}{$8a + 52$}
\solution{$(a+6)(a-6)-(a-4)^2 = (a^2-36)-(a^2-8a+16)=8a-52$. Відповідь: \textbf{А}.}
\zadtask{Спростіть вираз $(b+5)(b-5) - (b-3)^2$.}
\answerTable{$6b - 34$}{$-6b - 34$}{$6b - 16$}{$-34$}{$6b + 34$}
\zadtask{Спростіть вираз $(c-2)^2 - (c+7)(c-7)$.}
\answerTable{$-4c + 53$}{$4c - 53$}{$-4c - 45$}{$53$}{$-4c - 53$}
\taskBlock{Спрощення тригонометричного виразу --- сесія 5.06, завдання №10}
\zadtask{Спростіть вираз $\mathrm{tg}^2\alpha - \mathrm{tg}^2\alpha \cdot \sin^2\alpha$.}
\answerTableTall{$\sin^2\alpha$}{$\cos^2\alpha$}{$\mathrm{tg}^2\alpha$}{$1$}{$\sin^2\alpha\cos^2\alpha$}
\solution{$\mathrm{tg}^2\alpha-\mathrm{tg}^2\alpha\sin^2\alpha=\mathrm{tg}^2\alpha(1-\sin^2\alpha)=\dfrac{\sin^2\alpha}{\cos^2\alpha}\cdot\cos^2\alpha=\sin^2\alpha$. Відповідь: \textbf{А}.}
\zadtask{Спростіть вираз $\mathrm{ctg}^2\alpha - \mathrm{ctg}^2\alpha \cdot \cos^2\alpha$.}
\answerTableTall{$\sin^2\alpha$}{$\cos^2\alpha$}{$\mathrm{ctg}^2\alpha$}{$1$}{$\sin^2\alpha\cos^2\alpha$}
\zadtask{Спростіть вираз $\dfrac{1 - \sin^2\alpha}{\mathrm{ctg}^2\alpha}$.}
\answerTableTall{$\sin^2\alpha$}{$\cos^2\alpha$}{$\mathrm{ctg}^2\alpha$}{$1$}{$\mathrm{tg}^2\alpha$}
\typeTitle{Завдання на встановлення відповідності}
\taskBlock{Відповідність: вирази та їх значення --- сесія 23.05, завдання №17}
\zadtask{Установіть відповідність між виразом (1--3) та його значенням (А--Д). \nmtyear{2026}

\vspace{0.3cm}

\noindent
\begin{minipage}[t]{0.52\textwidth}
\textit{Вираз} \par \vspace{0.2cm}
\textbf{1} \ \ $\left|\dfrac{5}{3} - 2\right|$\\[0.35cm]
\textbf{2} \ \ $3^{-2} \cdot \left(\dfrac{1}{3}\right)^{-3}$\\[0.35cm]
\textbf{3} \ \ $\sqrt{3} \cdot \sin\dfrac{\pi}{3}$
\end{minipage}
\begin{minipage}[t]{0.24\textwidth}
\textit{Значення} \par \vspace{0.2cm}
\textbf{А} \ \ $\dfrac{3}{2}$\\[0.15cm]
\textbf{Б} \ \ $3$\\[0.15cm]
\textbf{В} \ \ $\dfrac{1}{3}$\\[0.15cm]
\textbf{Г} \ \ $-\dfrac{1}{3}$\\[0.15cm]
\textbf{Д} \ \ $\dfrac{1}{2}$
\end{minipage}
\hfill
\begin{minipage}[t]{0.22\textwidth}
\vspace{-0.2cm}
\begin{flushright}\matchingGrid\end{flushright}
\end{minipage}
}
\solution{$1)\ \left|\dfrac{5}{3}-2\right|=\left|-\dfrac{1}{3}\right|=\dfrac{1}{3}$~(В); $2)\ 3^{-2}\cdot 3^{3}=3^{1}=3$~(Б); $3)\ \sqrt{3}\cdot\sin\dfrac{\pi}{3}=\sqrt{3}\cdot\dfrac{\sqrt{3}}{2}=\dfrac{3}{2}$~(А). Відповідь: \textbf{1--В, 2--Б, 3--А}.}
\zadtask{Установіть відповідність між виразом (1--3) та його значенням (А--Д).

\vspace{0.3cm}

\noindent
\begin{minipage}[t]{0.52\textwidth}
\textit{Вираз} \par \vspace{0.2cm}
\textbf{1} \ \ $\left|\dfrac{7}{4} - 3\right|$\\[0.35cm]
\textbf{2} \ \ $2^{-3} \cdot \left(\dfrac{1}{2}\right)^{-5}$\\[0.35cm]
\textbf{3} \ \ $\sqrt{2} \cdot \cos\dfrac{\pi}{4}$
\end{minipage}
\begin{minipage}[t]{0.24\textwidth}
\textit{Значення} \par \vspace{0.2cm}
\textbf{А} \ \ $1$\\[0.15cm]
\textbf{Б} \ \ $\dfrac{5}{4}$\\[0.15cm]
\textbf{В} \ \ $\dfrac{1}{2}$\\[0.15cm]
\textbf{Г} \ \ $4$\\[0.15cm]
\textbf{Д} \ \ $-\dfrac{5}{4}$
\end{minipage}
\hfill
\begin{minipage}[t]{0.22\textwidth}
\vspace{-0.2cm}
\begin{flushright}\matchingGrid\end{flushright}
\end{minipage}
}
\zadtask{Установіть відповідність між виразом (1--3) та його значенням (А--Д).

\vspace{0.3cm}

\noindent
\begin{minipage}[t]{0.52\textwidth}
\textit{Вираз} \par \vspace{0.2cm}
\textbf{1} \ \ $\log_2 \dfrac{1}{4}$\\[0.35cm]
\textbf{2} \ \ $\left(\dfrac{2}{3}\right)^{-1}$\\[0.35cm]
\textbf{3} \ \ $2 \cdot \cos\dfrac{\pi}{3}$
\end{minipage}
\begin{minipage}[t]{0.24\textwidth}
\textit{Значення} \par \vspace{0.2cm}
\textbf{А} \ \ $1$\\[0.15cm]
\textbf{Б} \ \ $-2$\\[0.15cm]
\textbf{В} \ \ $\dfrac{3}{2}$\\[0.15cm]
\textbf{Г} \ \ $-\dfrac{1}{2}$\\[0.15cm]
\textbf{Д} \ \ $2$
\end{minipage}
\hfill
\begin{minipage}[t]{0.22\textwidth}
\vspace{-0.2cm}
\begin{flushright}\matchingGrid\end{flushright}
\end{minipage}
}
\taskBlock{Обчислення значень виразів (відповідність) --- сесія 30.05, завдання №18}
\zadtask{Установіть відповідність між виразом (1--3) і його значенням (А--Д), якщо $a = \dfrac{3}{8}$.

\vspace{0.2cm}
\noindent
\begin{minipage}[t]{0.45\textwidth}
\textit{Вираз}\par\vspace{0.15cm}
\textbf{1} \quad $\sqrt[3]{\dfrac{a}{3}}$\\[0.35cm]
\textbf{2} \quad $\log_2 a + \log_2 \dfrac{16}{3}$\\[0.35cm]
\textbf{3} \quad $\sin(4\pi a)$
\end{minipage}
\hfill
\begin{minipage}[t]{0.3\textwidth}
\textit{Значення}\par\vspace{0.15cm}
\textbf{А} \quad $2$\\[0.15cm]
\textbf{Б} \quad $1$\\[0.15cm]
\textbf{В} \quad $\dfrac{1}{2}$\\[0.2cm]
\textbf{Г} \quad $0$\\[0.15cm]
\textbf{Д} \quad $-1$
\end{minipage}
\hfill
\begin{minipage}[t]{0.2\textwidth}
\vspace{0.4cm}
\begin{flushright}\matchingGrid\end{flushright}
\end{minipage}
}
\solution{$a=\dfrac38$. \textbf{1}: $\sqrt[3]{\dfrac{a}{3}}=\sqrt[3]{\dfrac18}=\dfrac12$ (В). \textbf{2}: $\log_2\!\left(a\cdot\dfrac{16}{3}\right)=\log_2 2=1$ (Б). \textbf{3}: $\sin\!\left(4\pi\cdot\dfrac38\right)=\sin\dfrac{3\pi}{2}=-1$ (Д). Відповідь: \textbf{1~--~В, 2~--~Б, 3~--~Д}.}
\zadtask{Установіть відповідність між виразом (1--3) і його значенням (А--Д), якщо $a = \dfrac{5}{8}$.

\vspace{0.2cm}
\noindent
\begin{minipage}[t]{0.45\textwidth}
\textit{Вираз}\par\vspace{0.15cm}
\textbf{1} \quad $\sqrt[3]{\dfrac{a}{5}}$\\[0.35cm]
\textbf{2} \quad $\log_2 a + \log_2 \dfrac{32}{5}$\\[0.35cm]
\textbf{3} \quad $\sin(4\pi a)$
\end{minipage}
\hfill
\begin{minipage}[t]{0.3\textwidth}
\textit{Значення}\par\vspace{0.15cm}
\textbf{А} \quad $2$\\[0.15cm]
\textbf{Б} \quad $1$\\[0.15cm]
\textbf{В} \quad $\dfrac{1}{2}$\\[0.2cm]
\textbf{Г} \quad $0$\\[0.15cm]
\textbf{Д} \quad $-1$
\end{minipage}
\hfill
\begin{minipage}[t]{0.2\textwidth}
\vspace{0.4cm}
\begin{flushright}\matchingGrid\end{flushright}
\end{minipage}
}
\zadtask{Установіть відповідність між виразом (1--3) і його значенням (А--Д), якщо $a = \dfrac{1}{4}$.

\vspace{0.2cm}
\noindent
\begin{minipage}[t]{0.45\textwidth}
\textit{Вираз}\par\vspace{0.15cm}
\textbf{1} \quad $\sqrt{a}$\\[0.35cm]
\textbf{2} \quad $\log_2 \dfrac{1}{2a}$\\[0.35cm]
\textbf{3} \quad $\cos(4\pi a)$
\end{minipage}
\hfill
\begin{minipage}[t]{0.3\textwidth}
\textit{Значення}\par\vspace{0.15cm}
\textbf{А} \quad $2$\\[0.15cm]
\textbf{Б} \quad $1$\\[0.15cm]
\textbf{В} \quad $\dfrac{1}{2}$\\[0.2cm]
\textbf{Г} \quad $0$\\[0.15cm]
\textbf{Д} \quad $-1$
\end{minipage}
\hfill
\begin{minipage}[t]{0.2\textwidth}
\vspace{0.4cm}
\begin{flushright}\matchingGrid\end{flushright}
\end{minipage}
}
\taskBlock{Відповідність: значення виразів (степені, логарифми) --- сесія 1.06, завдання №16}
\zadtask{Узгодьте вираз (1--3) із його значенням (А--Д), якщо $a = 4$.

\vspace{0.3cm}
\noindent
\begin{minipage}[t]{0.45\textwidth}
\textit{Вираз}\par\vspace{0.2cm}
\textbf{1} \quad $\dfrac{2a^2 - 12a}{a^2 - 36}$\\[0.5cm]
\textbf{2} \quad $0{,}04^{\,\log_5 a}$\\[0.45cm]
\textbf{3} \quad $a^0 - a^{\frac{3}{2}}$
\end{minipage}
\hfill
\begin{minipage}[t]{0.3\textwidth}
\textit{Значення}\par\vspace{0.2cm}
\textbf{А} \quad $\dfrac{1}{16}$\\[0.25cm]
\textbf{Б} \quad $-8$\\[0.2cm]
\textbf{В} \quad $-7$\\[0.2cm]
\textbf{Г} \quad $0{,}8$\\[0.2cm]
\textbf{Д} \quad $-\dfrac{1}{16}$
\end{minipage}
\hfill
\begin{minipage}[t]{0.2\textwidth}
\vspace{0.5cm}
\begin{flushright}\matchingGrid\end{flushright}
\end{minipage}
}
\solution{При $a=4$: \textbf{1)} $\dfrac{2a^2-12a}{a^2-36}=\dfrac{2a(a-6)}{(a-6)(a+6)}=\dfrac{2a}{a+6}=\dfrac{8}{10}=0{,}8$ (Г); \textbf{2)} $0{,}04^{\log_5 4}=(5^{-2})^{\log_5 4}=4^{-2}=\dfrac{1}{16}$ (А); \textbf{3)} $a^0-a^{3/2}=1-8=-7$ (В). Відповідь: \textbf{1\,--\,Г, 2\,--\,А, 3\,--\,В}.}
\zadtask{Узгодьте вираз (1--3) із його значенням (А--Д), якщо $a = 9$.

\vspace{0.3cm}
\noindent
\begin{minipage}[t]{0.45\textwidth}
\textit{Вираз}\par\vspace{0.2cm}
\textbf{1} \quad $\dfrac{a^2 - 3a}{a^2 - 9}$\\[0.5cm]
\textbf{2} \quad $0{,}25^{\,\log_2 a}$\\[0.45cm]
\textbf{3} \quad $a^0 - a^{\frac{3}{2}}$
\end{minipage}
\hfill
\begin{minipage}[t]{0.3\textwidth}
\textit{Значення}\par\vspace{0.2cm}
\textbf{А} \quad $-26$\\[0.25cm]
\textbf{Б} \quad $\dfrac{1}{81}$\\[0.2cm]
\textbf{В} \quad $0{,}75$\\[0.2cm]
\textbf{Г} \quad $-\dfrac{1}{81}$\\[0.2cm]
\textbf{Д} \quad $-28$
\end{minipage}
\hfill
\begin{minipage}[t]{0.2\textwidth}
\vspace{0.5cm}
\begin{flushright}\matchingGrid\end{flushright}
\end{minipage}
}
\zadtask{Узгодьте вираз (1--3) із його значенням (А--Д), якщо $a = 16$.

\vspace{0.3cm}
\noindent
\begin{minipage}[t]{0.45\textwidth}
\textit{Вираз}\par\vspace{0.2cm}
\textbf{1} \quad $\dfrac{a - 9}{\sqrt{a} - 3}$\\[0.5cm]
\textbf{2} \quad $\log_2 a$\\[0.45cm]
\textbf{3} \quad $a^{-\frac{1}{2}}$
\end{minipage}
\hfill
\begin{minipage}[t]{0.3\textwidth}
\textit{Значення}\par\vspace{0.2cm}
\textbf{А} \quad $0{,}25$\\[0.25cm]
\textbf{Б} \quad $4$\\[0.2cm]
\textbf{В} \quad $7$\\[0.2cm]
\textbf{Г} \quad $16$\\[0.2cm]
\textbf{Д} \quad $2$
\end{minipage}
\hfill
\begin{minipage}[t]{0.2\textwidth}
\vspace{0.5cm}
\begin{flushright}\matchingGrid\end{flushright}
\end{minipage}
}
\taskBlock{Відповідність: значення виразу та проміжок --- сесія 2.06, завдання №16}
\zadtask{Узгодьте вираз (1--3) із проміжком (А--Д), якому належить значення цього виразу.

\vspace{0.3cm}
\noindent
\begin{minipage}[t]{0.4\textwidth}
\textit{Вираз}\par\vspace{0.2cm}
\textbf{1} \quad $\dfrac{7}{5} - \dfrac{5}{3}$\\[0.5cm]
\textbf{2} \quad $\sqrt{3} \cdot \sqrt{5}$\\[0.4cm]
\textbf{3} \quad $\log_4 8$
\end{minipage}
\hfill
\begin{minipage}[t]{0.32\textwidth}
\textit{Проміжок}\par\vspace{0.2cm}
\textbf{А} \quad $[-1;\, 0)$\\[0.15cm]
\textbf{Б} \quad $[0;\, 1)$\\[0.15cm]
\textbf{В} \quad $[1;\, 2)$\\[0.15cm]
\textbf{Г} \quad $[2;\, 4)$\\[0.15cm]
\textbf{Д} \quad $[4;\, +\infty)$
\end{minipage}
\hfill
\begin{minipage}[t]{0.2\textwidth}
\vspace{0.5cm}
\begin{flushright}\matchingGrid\end{flushright}
\end{minipage}
}
\solution{$\dfrac{7}{5}-\dfrac{5}{3}=\dfrac{21-25}{15}=-\dfrac{4}{15}\in[-1;0)$ -- \textbf{А}; $\sqrt{3}\cdot\sqrt{5}=\sqrt{15}\approx3{,}87\in[2;4)$ -- \textbf{Г}; $\log_4 8=\dfrac{3}{2}=1{,}5\in[1;2)$ -- \textbf{В}. Відповідь: \textbf{1--А, 2--Г, 3--В}.}
\zadtask{Узгодьте вираз (1--3) із проміжком (А--Д), якому належить значення цього виразу.

\vspace{0.3cm}
\noindent
\begin{minipage}[t]{0.4\textwidth}
\textit{Вираз}\par\vspace{0.2cm}
\textbf{1} \quad $\dfrac{9}{4} - \dfrac{5}{3}$\\[0.5cm]
\textbf{2} \quad $\sqrt{2} \cdot \sqrt{7}$\\[0.4cm]
\textbf{3} \quad $\log_9 27$
\end{minipage}
\hfill
\begin{minipage}[t]{0.32\textwidth}
\textit{Проміжок}\par\vspace{0.2cm}
\textbf{А} \quad $[-1;\, 0)$\\[0.15cm]
\textbf{Б} \quad $[0;\, 1)$\\[0.15cm]
\textbf{В} \quad $[1;\, 2)$\\[0.15cm]
\textbf{Г} \quad $[2;\, 4)$\\[0.15cm]
\textbf{Д} \quad $[4;\, +\infty)$
\end{minipage}
\hfill
\begin{minipage}[t]{0.2\textwidth}
\vspace{0.5cm}
\begin{flushright}\matchingGrid\end{flushright}
\end{minipage}
}
\zadtask{Узгодьте вираз (1--3) із проміжком (А--Д), якому належить значення цього виразу.

\vspace{0.3cm}
\noindent
\begin{minipage}[t]{0.4\textwidth}
\textit{Вираз}\par\vspace{0.2cm}
\textbf{1} \quad $\cos 60^\circ - 1$\\[0.5cm]
\textbf{2} \quad $\sqrt{50} - \sqrt{2}$\\[0.4cm]
\textbf{3} \quad $\log_2 12 - \log_2 3$
\end{minipage}
\hfill
\begin{minipage}[t]{0.32\textwidth}
\textit{Проміжок}\par\vspace{0.2cm}
\textbf{А} \quad $[-1;\, 0)$\\[0.15cm]
\textbf{Б} \quad $[0;\, 1)$\\[0.15cm]
\textbf{В} \quad $[1;\, 2)$\\[0.15cm]
\textbf{Г} \quad $[2;\, 4)$\\[0.15cm]
\textbf{Д} \quad $[4;\, +\infty)$
\end{minipage}
\hfill
\begin{minipage}[t]{0.2\textwidth}
\vspace{0.5cm}
\begin{flushright}\matchingGrid\end{flushright}
\end{minipage}
}
\taskBlock{Тотожні перетворення виразів (відповідність) --- сесія 3.06, завдання №16}
\noindent\zadnum Узгодьте вираз (1--3) з тотожно рівним йому виразом (А--Д), якщо $a \neq 0$.

\vspace{0.3cm}
\noindent
\begin{minipage}[t]{0.42\textwidth}
\textit{Вираз}\par\vspace{0.2cm}
\textbf{1} \quad $\sqrt[3]{\dfrac{27}{a^3}}$\\[0.55cm]
\textbf{2} \quad $a^{-2}\cdot a^3 \cdot \log_{27} 3$\\[0.4cm]
\textbf{3} \quad $\dfrac{(a-3)^2 - a^2 - 9}{2}$
\end{minipage}
\hfill
\begin{minipage}[t]{0.3\textwidth}
\textit{Тотожно рівний вираз}\par\vspace{0.2cm}
\textbf{А} \quad $-3a$\\[0.2cm]
\textbf{Б} \quad $\dfrac{3}{a}$\\[0.25cm]
\textbf{В} \quad $0$\\[0.2cm]
\textbf{Г} \quad $\dfrac{a}{3}$\\[0.25cm]
\textbf{Д} \quad $3a$
\end{minipage}
\hfill
\begin{minipage}[t]{0.2\textwidth}
\vspace{0.5cm}
\begin{flushright}\matchingGrid\end{flushright}
\end{minipage}

\solution{\textbf{1}: $\sqrt[3]{\dfrac{27}{a^3}} = \dfrac{3}{a}$~--- Б. \textbf{2}: $a^{-2}\cdot a^3\cdot\log_{27}3 = a\cdot\dfrac13 = \dfrac{a}{3}$~--- Г. \textbf{3}: $\dfrac{(a-3)^2 - a^2 - 9}{2} = \dfrac{-6a}{2} = -3a$~--- А. Відповідь: \textbf{1~--~Б, 2~--~Г, 3~--~А}.}

\vspace{0.4cm}
\noindent\zadnum Узгодьте вираз (1--3) з тотожно рівним йому виразом (А--Д), якщо $b \neq 0$.

\vspace{0.3cm}
\noindent
\begin{minipage}[t]{0.42\textwidth}
\textit{Вираз}\par\vspace{0.2cm}
\textbf{1} \quad $\sqrt[3]{\dfrac{64}{b^3}}$\\[0.55cm]
\textbf{2} \quad $b^{-3}\cdot b^4 \cdot \log_{16} 2$\\[0.4cm]
\textbf{3} \quad $\dfrac{(b-4)^2 - b^2 - 16}{2}$
\end{minipage}
\hfill
\begin{minipage}[t]{0.3\textwidth}
\textit{Тотожно рівний вираз}\par\vspace{0.2cm}
\textbf{А} \quad $\dfrac{4}{b}$\\[0.25cm]
\textbf{Б} \quad $4b$\\[0.2cm]
\textbf{В} \quad $\dfrac{b}{4}$\\[0.25cm]
\textbf{Г} \quad $-4b$\\[0.2cm]
\textbf{Д} \quad $0$
\end{minipage}
\hfill
\begin{minipage}[t]{0.2\textwidth}
\vspace{0.5cm}
\begin{flushright}\matchingGrid\end{flushright}
\end{minipage}

\vspace{0.4cm}
\noindent\zadnum Узгодьте вираз (1--3) з тотожно рівним йому виразом (А--Д), якщо $c \neq 0$.

\vspace{0.3cm}
\noindent
\begin{minipage}[t]{0.42\textwidth}
\textit{Вираз}\par\vspace{0.2cm}
\textbf{1} \quad $\left(\dfrac{c}{5}\right)^{-1}$\\[0.55cm]
\textbf{2} \quad $\sqrt[3]{-125c^3}$\\[0.4cm]
\textbf{3} \quad $\dfrac{(c+5)^2 - c^2 - 25}{2}$
\end{minipage}
\hfill
\begin{minipage}[t]{0.3\textwidth}
\textit{Тотожно рівний вираз}\par\vspace{0.2cm}
\textbf{А} \quad $5c$\\[0.2cm]
\textbf{Б} \quad $-5c$\\[0.2cm]
\textbf{В} \quad $\dfrac{5}{c}$\\[0.25cm]
\textbf{Г} \quad $0$\\[0.2cm]
\textbf{Д} \quad $\dfrac{c}{5}$
\end{minipage}
\hfill
\begin{minipage}[t]{0.2\textwidth}
\vspace{0.5cm}
\begin{flushright}\matchingGrid\end{flushright}
\end{minipage}
\taskBlock{Обчислення значень виразів (відповідність) --- сесія 4.06, завдання №16}
\zadtask{Узгодьте вираз (1--3) із його значенням (А--Д), якщо $a = \dfrac{2}{3}$.

\vspace{0.3cm}
\noindent
\begin{minipage}[t]{0.4\textwidth}
\textit{Вираз}\par\vspace{0.2cm}
\textbf{1} \quad $6a$\\[0.4cm]
\textbf{2} \quad $\sqrt{\left(2^a\right)^3}$\\[0.4cm]
\textbf{3} \quad $\log_3 (1 - a)$
\end{minipage}
\hfill
\begin{minipage}[t]{0.32\textwidth}
\textit{Значення виразу}\par\vspace{0.2cm}
\textbf{А} \quad $4$\\[0.25cm]
\textbf{Б} \quad $-1$\\[0.25cm]
\textbf{В} \quad $1$\\[0.25cm]
\textbf{Г} \quad $\dfrac{20}{3}$\\[0.25cm]
\textbf{Д} \quad $2$
\end{minipage}
\hfill
\begin{minipage}[t]{0.2\textwidth}
\vspace{0.4cm}
\begin{flushright}\matchingGrid\end{flushright}
\end{minipage}
}
\solution{При $a=\dfrac{2}{3}$: $6a=4$ (А); $\sqrt{(2^a)^3}=\sqrt{2^{3a}}=\sqrt{2^2}=2$ (Д); $\log_3(1-a)=\log_3\dfrac{1}{3}=-1$ (Б). Відповідь: \textbf{1~--~А, 2~--~Д, 3~--~Б}.}
\zadtask{Узгодьте вираз (1--3) із його значенням (А--Д), якщо $a = \dfrac{3}{2}$.

\vspace{0.3cm}
\noindent
\begin{minipage}[t]{0.4\textwidth}
\textit{Вираз}\par\vspace{0.2cm}
\textbf{1} \quad $4a$\\[0.4cm]
\textbf{2} \quad $\sqrt{2^{4a}}$\\[0.4cm]
\textbf{3} \quad $\log_2 \left(a + \dfrac{1}{2}\right)$
\end{minipage}
\hfill
\begin{minipage}[t]{0.32\textwidth}
\textit{Значення виразу}\par\vspace{0.2cm}
\textbf{А} \quad $1$\\[0.25cm]
\textbf{Б} \quad $6$\\[0.25cm]
\textbf{В} \quad $8$\\[0.25cm]
\textbf{Г} \quad $3$\\[0.25cm]
\textbf{Д} \quad $-1$
\end{minipage}
\hfill
\begin{minipage}[t]{0.2\textwidth}
\vspace{0.4cm}
\begin{flushright}\matchingGrid\end{flushright}
\end{minipage}
}
\zadtask{Узгодьте вираз (1--3) із його значенням (А--Д), якщо $a = -2$.

\vspace{0.3cm}
\noindent
\begin{minipage}[t]{0.4\textwidth}
\textit{Вираз}\par\vspace{0.2cm}
\textbf{1} \quad $a^3$\\[0.4cm]
\textbf{2} \quad $|a| + a^2$\\[0.4cm]
\textbf{3} \quad $\left(\dfrac{1}{2}\right)^{a}$
\end{minipage}
\hfill
\begin{minipage}[t]{0.32\textwidth}
\textit{Значення виразу}\par\vspace{0.2cm}
\textbf{А} \quad $6$\\[0.25cm]
\textbf{Б} \quad $-8$\\[0.25cm]
\textbf{В} \quad $4$\\[0.25cm]
\textbf{Г} \quad $-2$\\[0.25cm]
\textbf{Д} \quad $2$
\end{minipage}
\hfill
\begin{minipage}[t]{0.2\textwidth}
\vspace{0.4cm}
\begin{flushright}\matchingGrid\end{flushright}
\end{minipage}
}
\taskBlock{Обчислення значень виразів (відповідність) --- сесія 5.06, завдання №16}
\zadtask{Узгодьте вираз (1--3) із його значенням (А--Д).

\vspace{0.3cm}
\noindent
\begin{minipage}[t]{0.4\textwidth}
\textit{Вираз}\par\vspace{0.2cm}
\textbf{1} \quad $\left|\,-6 + \cos 4\pi\,\right|$\\[0.4cm]
\textbf{2} \quad $\sqrt{2}\cdot\sqrt{18}$\\[0.4cm]
\textbf{3} \quad $\log_{1/2} 32$
\end{minipage}
\hfill
\begin{minipage}[t]{0.3\textwidth}
\textit{Значення}\par\vspace{0.2cm}
\textbf{А} \quad $-6$\\[0.15cm]
\textbf{Б} \quad $-5$\\[0.15cm]
\textbf{В} \quad $5$\\[0.15cm]
\textbf{Г} \quad $6$\\[0.15cm]
\textbf{Д} \quad $7$
\end{minipage}
\hfill
\begin{minipage}[t]{0.2\textwidth}
\vspace{0.5cm}
\begin{flushright}\matchingGrid\end{flushright}
\end{minipage}}
\solution{$1$: $|-6+\cos4\pi|=|-6+1|=5$ (В); $2$: $\sqrt{2}\cdot\sqrt{18}=\sqrt{36}=6$ (Г); $3$: $\log_{1/2}32=\log_{1/2}2^5=-5$ (Б). Відповідь: \textbf{1~--~В, 2~--~Г, 3~--~Б}.}
\zadtask{Узгодьте вираз (1--3) із його значенням (А--Д).

\vspace{0.3cm}
\noindent
\begin{minipage}[t]{0.4\textwidth}
\textit{Вираз}\par\vspace{0.2cm}
\textbf{1} \quad $\left|\,-8 + \cos 2\pi\,\right|$\\[0.4cm]
\textbf{2} \quad $\sqrt{3}\cdot\sqrt{27}$\\[0.4cm]
\textbf{3} \quad $\log_{1/3} 81$
\end{minipage}
\hfill
\begin{minipage}[t]{0.3\textwidth}
\textit{Значення}\par\vspace{0.2cm}
\textbf{А} \quad $-7$\\[0.15cm]
\textbf{Б} \quad $-4$\\[0.15cm]
\textbf{В} \quad $7$\\[0.15cm]
\textbf{Г} \quad $9$\\[0.15cm]
\textbf{Д} \quad $-9$
\end{minipage}
\hfill
\begin{minipage}[t]{0.2\textwidth}
\vspace{0.5cm}
\begin{flushright}\matchingGrid\end{flushright}
\end{minipage}}
\zadtask{Узгодьте вираз (1--3) із його значенням (А--Д).

\vspace{0.3cm}
\noindent
\begin{minipage}[t]{0.4\textwidth}
\textit{Вираз}\par\vspace{0.2cm}
\textbf{1} \quad $2^{-3} \cdot 2^{5}$\\[0.4cm]
\textbf{2} \quad $\sqrt[3]{-64}$\\[0.4cm]
\textbf{3} \quad $6\cos\pi$
\end{minipage}
\hfill
\begin{minipage}[t]{0.3\textwidth}
\textit{Значення}\par\vspace{0.2cm}
\textbf{А} \quad $-6$\\[0.15cm]
\textbf{Б} \quad $-4$\\[0.15cm]
\textbf{В} \quad $4$\\[0.15cm]
\textbf{Г} \quad $10$\\[0.15cm]
\textbf{Д} \quad $2$
\end{minipage}
\hfill
\begin{minipage}[t]{0.2\textwidth}
\vspace{0.5cm}
\begin{flushright}\matchingGrid\end{flushright}
\end{minipage}}
\typeTitle{Завдання з короткою відповіддю}
\taskBlock{Сюжетна задача на відсотки і частини --- сесія 23.05, завдання №21}
\zadtask{У магазині є лише музичні диски, диски з науково-популярними фільмами та диски з художніми фільмами. Кількість музичних дисків утричі більша за кількість дисків із науково-популярними фільмами й становить $40\%$ від кількості дисків з художніми фільмами. Загальна кількість дисків $345$. Визначте кількість музичних дисків у магазині. \nmtyear{2026}}
\nmtAnswerBox
\solution{Нехай науково-популярних~--- $x$, тоді музичних~--- $3x$, а художніх~--- $\dfrac{3x}{0{,}4}=7{,}5x$. Сума: $x+3x+7{,}5x=11{,}5x=345$, звідки $x=30$. Музичних: $3x=90$. Відповідь: \textbf{90}.}
\zadtask{У бібліотеці є лише довідники, словники та енциклопедії. Кількість довідників удвічі більша за кількість словників і становить $25\%$ від кількості енциклопедій. Загальна кількість книжок $220$. Визначте кількість довідників у бібліотеці.}
\nmtAnswerBox
\zadtask{У кошику є лише яблука, груші та сливи. Кількість яблук удвічі менша за кількість груш і становить $20\%$ від кількості слив. У кошику $12$ яблук. Визначте загальну кількість фруктів у кошику.}
\nmtAnswerBox
\taskBlock{Задача на відсотки (суміш) --- сесія 30.05, завдання №20}
\zadtask{Для приготування трав'яного чаю змішали м'яту й чебрець. Маса суміші становить $208$ г, причому м'яти в суміші на $60\,\%$ більше, ніж чебрецю. Скільки грамів м'яти у цій суміші?
\nmtAnswerBox}
\solution{Нехай чебрецю $c$ г, тоді м'яти $1{,}6c$ г. Маємо $c+1{,}6c=208$, звідки $2{,}6c=208$, $c=80$. М'яти: $1{,}6\cdot80=128$ г. Відповідь: $128$.}
\zadtask{Для приготування трав'яного збору змішали ромашку й липу. Маса збору становить $252$ г, причому ромашки у зборі на $80\,\%$ більше, ніж липи. Скільки грамів ромашки у цьому зборі?
\nmtAnswerBox}
\zadtask{Для приготування мюслі змішали вівсяні пластівці й горіхи. Маса суміші становить $170$ г, причому горіхів у суміші на $30\,\%$ менше, ніж пластівців. Скільки грамів горіхів у цій суміші?
\nmtAnswerBox}
\taskBlock{Текстова задача на відсотки --- сесія 1.06, завдання №20}
\zadtask{У будинку $162$ квартири~--- одно-, дво- і трикімнатні. Двокімнатних квартир на $15\,\%$ менше, ніж однокімнатних, а трикімнатних~--- стільки ж, скільки двокімнатних. Скільки в будинку \textit{двокімнатних} квартир?}
\nmtAnswerBox
\solution{Нехай однокімнатних $x$, тоді дво- і трикімнатних по $0{,}85x$. Усього $x+0{,}85x+0{,}85x=2{,}7x=162$, звідки $x=60$. Двокімнатних: $0{,}85\cdot 60=51$. Відповідь: \textbf{$51$}.}
\zadtask{У будинку $150$ квартир~--- одно-, дво- і трикімнатні. Двокімнатних квартир на $25\,\%$ менше, ніж однокімнатних, а трикімнатних~--- стільки ж, скільки двокімнатних. Скільки в будинку \textit{двокімнатних} квартир?}
\nmtAnswerBox
\zadtask{У гаражному кооперативі $240$ боксів~--- малі, середні та великі. Малих боксів на $20\,\%$ \textit{більше}, ніж великих, а середніх~--- стільки ж, скільки великих. Скільки в кооперативі \textit{великих} боксів?}
\nmtAnswerBox
\taskBlock{Сюжетна задача на відсотки (вартість таксі) --- сесія 2.06, завдання №20}
\zadtask{Перед подорожжю турист проаналізував інфографіку й дізнався вартість (у грн) поїздки від аеропорту до центру міста (або у зворотному напрямку) у деяких світових столицях. На час туристичного сезону вартість проїзду в таксі в Осло збільшується на $20\,\%$, а в Рейк'явіку~--- на $30\,\%$ від фіксованої вартості проїзду. Визначте сумарну кількість коштів, яку витратить турист на таксі від аеропорту до центру міста й назад у цих містах, якщо відвідає Осло та Рейк'явік у туристичний сезон.

\vspace{0.3cm}
\begin{center}
\begin{tikzpicture}[scale=1.0]
    % фон-рамка
    \draw[gray!40, rounded corners=3pt] (-0.3,-0.4) rectangle (9.2,5.6);
    % дані: місто / значення / довжина бара (масштаб: 5561 -> 4.2)
    \def\scale{0.000755}
    \foreach \i/\city/\val in {0/{аеропорт Лондона}/2502, 1/{аеропорт Монако}/2502, 2/{аеропорт Рейк'явіка}/2780, 3/{аеропорт Осло}/2780, 4/{аеропорт Токіо}/5561} {
        \pgfmathsetmacro\y{\i*1.05}
        \pgfmathsetmacro\barlen{\val*\scale}
        % прапорець-заглушка
        \filldraw[fill=gray!25, draw=gray!50] (0,\y+0.32) rectangle (0.45,\y+0.62);
        % назва
        \node[right, font=\small] at (0.55,\y+0.47) {\city};
        % бар
        \filldraw[fill=cyan!45, draw=cyan!60!black] (0.55,\y) rectangle (0.55+\barlen,\y+0.28);
        \node[right, font=\scriptsize\bfseries, text=white] at (0.7,\y+0.14) {\val};
    }
    % легенда
    \filldraw[fill=cyan!45, draw=cyan!60!black] (6.0,4.55) rectangle (6.35,4.85);
    \node[right, font=\scriptsize] at (6.4,4.7) {Фіксована вартість};
    \node[right, font=\scriptsize] at (6.4,4.4) {проїзду, \textit{грн}};
\end{tikzpicture}
\end{center}
\nmtAnswerBox
}
\solution{Фіксована вартість в Осло й Рейк'явіку по $2780$ грн. У сезон: Осло $2780\cdot1{,}2$, Рейк'явік $2780\cdot1{,}3$, а туди й назад -- подвоюємо: $2\cdot(2780\cdot1{,}2+2780\cdot1{,}3)=2\cdot2780\cdot2{,}5=13900$ грн.}
\zadtask{Перед подорожжю турист проаналізував інфографіку й дізнався вартість (у грн) поїздки від аеропорту до центру міста (або у зворотному напрямку) у деяких європейських столицях. На час туристичного сезону вартість проїзду в таксі в Парижі збільшується на $25\,\%$, а в Римі~--- на $10\,\%$ від фіксованої вартості проїзду. Визначте сумарну кількість коштів, яку витратить турист на таксі від аеропорту до центру міста й назад у цих містах, якщо відвідає Париж та Рим у туристичний сезон.

\vspace{0.3cm}
\begin{center}
\begin{tikzpicture}[scale=1.0]
    % фон-рамка
    \draw[gray!40, rounded corners=3pt] (-0.3,-0.4) rectangle (9.2,5.6);
    \def\scale{0.0016}
    \foreach \i/\city/\val in {0/{аеропорт Праги}/1500, 1/{аеропорт Відня}/1800, 2/{аеропорт Мадрида}/2000, 3/{аеропорт Парижа}/2400, 4/{аеропорт Рима}/2600} {
        \pgfmathsetmacro\y{\i*1.05}
        \pgfmathsetmacro\barlen{\val*\scale}
        \filldraw[fill=gray!25, draw=gray!50] (0,\y+0.32) rectangle (0.45,\y+0.62);
        \node[right, font=\small] at (0.55,\y+0.47) {\city};
        \filldraw[fill=cyan!45, draw=cyan!60!black] (0.55,\y) rectangle (0.55+\barlen,\y+0.28);
        \node[right, font=\scriptsize\bfseries, text=white] at (0.7,\y+0.14) {\val};
    }
    \filldraw[fill=cyan!45, draw=cyan!60!black] (6.0,4.55) rectangle (6.35,4.85);
    \node[right, font=\scriptsize] at (6.4,4.7) {Фіксована вартість};
    \node[right, font=\scriptsize] at (6.4,4.4) {проїзду, \textit{грн}};
\end{tikzpicture}
\end{center}
\nmtAnswerBox
}
\zadtask{Перед подорожжю турист проаналізував інфографіку й дізнався вартість (у грн) поїздки від аеропорту до центру міста (або у зворотному напрямку) у деяких європейських столицях. У міжсезоння вартість проїзду в таксі у Відні \textit{зменшується} на $10\,\%$, а в Мадриді~--- на $20\,\%$ від фіксованої вартості проїзду. Визначте сумарну кількість коштів, яку витратить турист на таксі від аеропорту до центру міста й назад у цих містах, якщо відвідає Відень та Мадрид у міжсезоння.

\vspace{0.3cm}
\begin{center}
\begin{tikzpicture}[scale=1.0]
    % фон-рамка
    \draw[gray!40, rounded corners=3pt] (-0.3,-0.4) rectangle (9.2,5.6);
    \def\scale{0.0016}
    \foreach \i/\city/\val in {0/{аеропорт Праги}/1500, 1/{аеропорт Відня}/1800, 2/{аеропорт Мадрида}/2000, 3/{аеропорт Парижа}/2400, 4/{аеропорт Рима}/2600} {
        \pgfmathsetmacro\y{\i*1.05}
        \pgfmathsetmacro\barlen{\val*\scale}
        \filldraw[fill=gray!25, draw=gray!50] (0,\y+0.32) rectangle (0.45,\y+0.62);
        \node[right, font=\small] at (0.55,\y+0.47) {\city};
        \filldraw[fill=cyan!45, draw=cyan!60!black] (0.55,\y) rectangle (0.55+\barlen,\y+0.28);
        \node[right, font=\scriptsize\bfseries, text=white] at (0.7,\y+0.14) {\val};
    }
    \filldraw[fill=cyan!45, draw=cyan!60!black] (6.0,4.55) rectangle (6.35,4.85);
    \node[right, font=\scriptsize] at (6.4,4.7) {Фіксована вартість};
    \node[right, font=\scriptsize] at (6.4,4.4) {проїзду, \textit{грн}};
\end{tikzpicture}
\end{center}
\nmtAnswerBox
}
\taskBlock{Відсотки і знижки (сюжетна задача) --- сесія 3.06, завдання №20}
\zadtask{Сім'я з трьох осіб~--- батька, мами й сина~--- запланувала сплав на байдарці. Вартість сплаву~--- $800$ грн з особи. Мамі та сину надають знижки: $20\,\%$ і $25\,\%$ відповідно (батько сплачує повну вартість). Скільки гривень заплатить сім'я за сплав?}
\nmtAnswerBox
\solution{Батько платить повну вартість $800$ грн, мама $800\cdot0{,}8 = 640$ грн, син $800\cdot0{,}75 = 600$ грн. Усього $800 + 640 + 600 = 2040$ грн. Відповідь: \textbf{$2040$}.}
\zadtask{Сім'я з трьох осіб~--- батька, мами й доньки~--- запланувала екскурсію до печер. Вартість екскурсії~--- $600$ грн з особи. Мамі та доньці надають знижки: $10\,\%$ і $50\,\%$ відповідно (батько сплачує повну вартість). Скільки гривень заплатить сім'я за екскурсію?}
\nmtAnswerBox
\zadtask{Сім'я з трьох осіб~--- батька, мами й сина~--- купує квитки на виставу. Вартість квитка~--- $900$ грн. Мамі та сину надають знижки: $20\,\%$ і $30\,\%$ відповідно (батько сплачує повну вартість квитка). \textit{На скільки гривень менше} заплатить сім'я порівняно з повною вартістю трьох квитків?}
\nmtAnswerBox
\taskBlock{Текстова задача на відсотки --- сесія 4.06, завдання №19}
\zadtask{Валерія за два дні прочитала $132$ вірші. За перший день вона прочитала на $20\,\%$ віршів більше, ніж за другий. Скільки віршів прочитала Валерія за перший день?
\nmtAnswerBox}
\solution{Нехай за другий день прочитано $x$ віршів, тоді за перший $1{,}2x$. Сума: $x+1{,}2x=132$, $2{,}2x=132$, $x=60$. За перший день $1{,}2\cdot60=72$. Відповідь: \textbf{72}.}
\zadtask{Остап за два дні прочитав $180$ сторінок книжки. За перший день він прочитав на $25\,\%$ сторінок більше, ніж за другий. Скільки сторінок прочитав Остап за перший день?
\nmtAnswerBox}
\zadtask{Соломія за два дні розв'язала $119$ задач. За перший день вона розв'язала на $30\,\%$ задач менше, ніж за другий. Скільки задач розв'язала Соломія за перший день?
\nmtAnswerBox}
\taskBlock{Сюжетна задача на відсотки (акційна покупка) --- сесія 5.06, завдання №19}
\zadtask{Чоловік купив три футболки за $1056$ грн. У магазині діяла акція: першу футболку продавали за повну ціну, другу~--- на $20\,\%$ дешевше, а третю~--- на $40\,\%$ дешевше за повну ціну. Скільки гривень коштувала одна футболка \textit{без} акції?}
\nmtAnswerBox
\solution{Нехай повна ціна футболки $x$ грн. Тоді $x+0{,}8x+0{,}6x=2{,}4x=1056$, звідки $x=440$. Відповідь: \textbf{$440$}.}
\zadtask{Жінка купила три однакові чашки за $891$ грн. У магазині діяла акція: першу чашку продавали за повну ціну, другу~--- на $30\,\%$ дешевше, а третю~--- на $50\,\%$ дешевше за повну ціну. Скільки гривень коштувала одна чашка \textit{без} акції?}
\nmtAnswerBox
\zadtask{Покупець придбав три однакові рюкзаки й заплатив за них $1840$ грн. У магазині діяла акція: перший рюкзак продавали за повну ціну, другий~--- на $25\,\%$ дешевше, а третій~--- на $45\,\%$ дешевше за повну ціну. Скільки гривень \textit{заощадив} покупець завдяки акції порівняно з купівлею трьох рюкзаків за повною ціною?}
\nmtAnswerBox
\chapterTitle{ВІДПОВІДІ}
\noindent\small Відповіді наведено у тому ж порядку, що й завдання.\par\vspace{0.3cm}
\ansType{Завдання з вибором однієї відповіді (А--Д)}
\begin{multicols}{3}\noindent
\ansitem{А}
\ansitem{Б}
\ansitem{Б}
\ansitem{В}
\ansitem{А}
\ansitem{А}
\ansitem{В}
\ansitem{Б}
\ansitem{Б}
\ansitem{А}
\ansitem{Б}
\ansitem{А}
\ansitem{Б}
\ansitem{А}
\ansitem{Б}
\ansitem{Г}
\ansitem{В}
\ansitem{А}
\ansitem{Г}
\ansitem{В}
\ansitem{В}
\ansitem{Д}
\ansitem{Д}
\ansitem{Д}
\ansitem{А}
\ansitem{Б}
\ansitem{А}
\ansitem{А}
\ansitem{Б}
\ansitem{В}
\ansitem{А}
\ansitem{А}
\ansitem{Г}
\ansitem{В}
\ansitem{Б}
\ansitem{Б}
\ansitem{А}
\ansitem{Б}
\ansitem{Б}
\ansitem{Б}
\ansitem{Б}
\ansitem{А}
\ansitem{А}
\ansitem{В}
\ansitem{В}
\ansitem{Б}
\ansitem{В}
\ansitem{Д}
\ansitem{Г}
\ansitem{В}
\ansitem{Б}
\ansitem{А}
\ansitem{А}
\ansitem{А}
\ansitem{А}
\ansitem{Б}
\ansitem{А}
\end{multicols}
\ansType{Завдання на встановлення відповідності}
\begin{multicols}{3}\noindent
\ansitem{1~--~В, 2~--~Б, 3~--~А}
\ansitem{1~--~Б, 2~--~Г, 3~--~А}
\ansitem{1~--~Б, 2~--~В, 3~--~А}
\ansitem{1~--~В, 2~--~Б, 3~--~Д}
\ansitem{1~--~В, 2~--~А, 3~--~Б}
\ansitem{1~--~В, 2~--~Б, 3~--~Д}
\ansitem{1~--~Г, 2~--~А, 3~--~В}
\ansitem{1~--~В, 2~--~Б, 3~--~А}
\ansitem{1~--~В, 2~--~Б, 3~--~А}
\ansitem{1~--~А, 2~--~Г, 3~--~В}
\ansitem{1~--~Б, 2~--~Г, 3~--~В}
\ansitem{1~--~А, 2~--~Д, 3~--~Г}
\ansitem{1~--~Б, 2~--~Г, 3~--~А}
\ansitem{1~--~А, 2~--~В, 3~--~Г}
\ansitem{1~--~В, 2~--~Б, 3~--~А}
\ansitem{1~--~А, 2~--~Д, 3~--~Б}
\ansitem{1~--~Б, 2~--~В, 3~--~А}
\ansitem{1~--~Б, 2~--~А, 3~--~В}
\ansitem{1~--~В, 2~--~Г, 3~--~Б}
\ansitem{1~--~В, 2~--~Г, 3~--~Б}
\ansitem{1~--~В, 2~--~Б, 3~--~А}
\end{multicols}
\ansType{Завдання з короткою відповіддю}
\begin{multicols}{3}\noindent
\ansitem{$90$}
\ansitem{$40$}
\ansitem{$96$}
\ansitem{$128$}
\ansitem{$162$}
\ansitem{$70$}
\ansitem{$51$}
\ansitem{$45$}
\ansitem{$75$}
\ansitem{$13900$}
\ansitem{$11720$}
\ansitem{$6440$}
\ansitem{$2040$}
\ansitem{$1440$}
\ansitem{$450$}
\ansitem{$72$}
\ansitem{$100$}
\ansitem{$49$}
\ansitem{$440$}
\ansitem{$405$}
\ansitem{$560$}
\end{multicols}

\end{document}
